%% kapa: exam=1 round=26 year=2015 period=2 form=A subject=appraisal_law subject_ko=부동산관계법규 seq=39 number=39
%% source: Q-Net 감정평가사 기출문제 (cst00309) artlSeq=5138001 file=extracted/1008/5138001/제26회_감정평가사 1차_기출/제26회_감정평가사 1차_기출/2교시_감평_A형.pdf
%% conversion: kordoc 4.24.51 -> tools/gen1.py (auto); review=proofread-sonnet; flags=none
\begin{question}{39}
부동산등기법령상 등기관이 직권으로 등기를 말소할 수 있는 경우는?
\begin{choices}
  \item 등기를 신청할 때 당사자가 출석하지 아니한 경우
  \item 등기에 필요한 첨부정보를 제공하지 아니한 경우
  \item 신청할 권한이 없는 자가 신청한 경우
  \item 사건이 등기할 것이 아닌 경우
  \item 신청정보와 등기원인을 증명하는 정보가 일치하지 아니한 경우
\end{choices}
\end{question}
\answer{4}
